\section{Local geometric statements and conventions}
\label{sec:introduction84}
This supplement keeps the local part of the ordered-measurement theory
separate from the exact finite decision problem.  We first recall its
principal statements and notation.  The labelled summary formulas also
preserve the references used by the retained proofs.

An ordered measurement $E=(E_1,\ldots,E_k)$ consists of positive
$d\times d$ matrices summing to $I$.  Put
$\|R\|_\Sigma=\sum_j\|R_j\|_{\rm op}$.  Fix $E$, a nonzero
Hermitian zero-sum one-sided tangent $H$, and $\Lambda\ge0$.  The
legal quadratic neighborhood is
\begin{equation}\label{eq:fronttube87}
                 \|F-E-sH\|_\Sigma\le\Lambda s^2.
\end{equation}
For $P_j=\operatorname{supp}E_j$, $Q_j=I-P_j$, put
$J_j=Q_jH_jQ_j$, $\Gamma_E(H)=\frac i2\sum_j[P_j,H_j]$, and
$\mathcal W_E=\sum_jP_j\mathcal H_dP_j$.  The cone condition is
$J_j\succeq0$.  For positive integer complete acquisition sizes
$n_1,\ldots,n_\ell$, write $N=\sum_rn_r$, $V=\sum_rn_r^2$.
Theorem~\ref{thm:allocation89} gives the respective local orders
\begin{equation}\label{eq:frontprofile89}
 \min\{1,\sqrt N s\},\qquad\min\{1,\sqrt V s\},\qquad\min\{1,Ns\}
\end{equation}
for $H\in\operatorname{Ran}\mathcal C_E$, for all $J_j=0$ with
$\Gamma_E(H)\notin\mathcal W_E$, and for some $J_j\ne0$.
The support-kernel identity exhausts the nonzero cone.  The constants
and the small-scale interval depend on the fixed $E,H,\Lambda$,
not on the allocation or the unknown remainder.  Prescribed-profile
reset feedback has the same orders only with conditional independence
of the entire fresh acquisition from the entire old receiver.

For a width cap $1\le b\le N$ the consequence is
\begin{equation}\label{eq:frontwidth87}
 D_N^{[b]}(E,F)\asymp_{E,H,\Lambda}
 \begin{cases}
 \min\{1,\sqrt Ns\},&H\in\operatorname{Ran}\mathcal C_E,\\
 \min\{1,\sqrt{Nb}\,s\},&J_j=0\ (\forall j),\ \Gamma_E(H)\notin\mathcal W_E,\\
 \min\{1,Ns\},&J_j\ne0\ (\exists j).
 \end{cases}
\end{equation}
Hard policy costs use suprema over all formal paths, not expected
costs.  Theorem~\ref{thm:quadraticbudget89} uses
$Q_\pi=\sup_\zeta\sum_{h\in\zeta}n(h)^2$ and gives coherent
order $\min\{1,\sqrt Qs\}$ at its stated hard budgets.

With $M=(E+F)/2$ and $H=F-E$, the covariance upper certificate is
\begin{equation}\label{eq:frontupper84}
 D_N(E,F)\le\min\left\{2,2\sqrt{k+1}
 \sqrt{N\langle H,(\mathcal C_M+N^{-1}\operatorname{Id})^{-1}H\rangle}
 \right\}.
\end{equation}
This is not a global matching lower or uniform full-boundary metric.
On a fixed nonstationary orbit $E_j(t)=e^{itG}E_je^{-itG}$,
Theorem~\ref{thm:orbittrichotomy84} gives
\begin{equation}\label{eq:frontorbit84}
 D_N(E,E(t))\asymp_{E,G}
 \begin{cases}
 \min\{1,\sqrt N|t|\},&G\in\mathcal W_E,\\
 \min\{1,N|t|\},&G\notin\mathcal W_E.
 \end{cases}
\end{equation}
The finite remainder estimates, support-span comparison and physical
controls are proved in the following sections with their original
hypotheses.

For comparison only, the independent exact two-call shared-basis
experiment of the primary uses $L=2(d+1)-t^2$ and has
\begin{equation}\label{eq:fronthierarchy90}
 \begin{split}
 P_{\rm ca}&=(d+1)/L,\\
 P_{\rm reset}&=(d+1)/L+t^2(d-1)/(2dL),\\
 P_{\rm all}&=(d+1+t^2)/L.
 \end{split}
\end{equation}
Its once-selected alternative and prescribed priors are those of
Theorem~\ref{thm:operationalhierarchy90}; it is not a consequence of
\eqref{eq:frontprofile89}.  All trace norms in the local results are
unhalved.  Reference systems, complete reset cuts and ordinary Bayes
probabilities retain the primary's explicit conventions.
